\chapter{相关技术与原理}

本章介绍本系统所依赖的核心理论与关键技术，包括人眼视觉暂留与视觉积分特性、时间分片与高刷新率显示、密码学安全伪随机数生成、对抗样本技术以及光学字符识别与目标检测技术，为后续章节的系统设计提供理论支撑。

\section{视觉暂留与人眼视觉积分}

视觉暂留是指光刺激在视网膜上消失后，视觉感知仍能短暂保留一段时间的生理现象。人眼的感光细胞在受到光照后，其神经信号并不会随光照消失而立即归零，而是经历一个约 50ms 至 100ms 的衰减过程。这一特性使得人眼在感知快速变化的光信号时，等效于对一段时间窗口内的光强进行积分，可以近似建模为一个低通滤波器。电影、电视等动态影像技术正是利用这一特性，以每秒 24 帧以上的速度连续播放静态画面，使人眼感知到连续运动。

设第 $k$ 个子帧在时刻 $t_k$ 显示、持续时间为 $\Delta t$，子帧光强为 $L_k(x,y)$，则人眼在积分窗口 $T$ 内于像素 $(x,y)$ 处感知的有效光强可表示为：

\begin{equation}
E(x,y) = \frac{1}{n}\sum_{k=1}^{n} L_k(x,y) \cdot \Delta t \cdot f_r
\end{equation}

其中 $f_r$ 为刷新率，$n$ 为一个完整周期内的子帧数。当完整周期时长 $T_{cycle}=n/f_r \le 50$ms 时，人眼无法分辨单个子帧，只能感知到所有子帧叠加后的积分结果。本系统正是利用人眼这种"时域积分"特性与相机"瞬时快门采样"特性之间的根本差异，实现"人眼可见、相机不可识别"的效果。

\section{时间分片与高刷新率显示}

时间分片是指将一帧完整图像在时间维度上拆分为多个子帧，分时轮流显示的技术。为保证人眼无可感知闪烁，一个完整周期必须落在临界融合频率对应的时间窗口内。临界闪烁融合频率（Critical Flicker Fusion, CFF）是人眼能够分辨光强闪烁的最高频率，一般在 60Hz 左右。因此，要使 $n$ 个子帧组成的完整周期不被人眼察觉，需满足：

\begin{equation}
T_{cycle} = \frac{n}{f_r} \le \frac{1}{60}\,\text{s} \approx 16.7\,\text{ms},\quad\text{即}\quad f_r \ge 60n
\end{equation}

由此可得：$n=4$ 时需 $f_r \ge 240$Hz，$n=6$ 时需 $f_r \ge 360$Hz，$n=8$ 时需 $f_r \ge 480$Hz。近年来电竞显示器刷新率已普遍达到 240Hz 乃至 480Hz，为时间分片方案提供了硬件基础。需要说明的是，本系统作为概念验证原型，在普通 60Hz 至 120Hz 显示器上以缩小的 $n$ 值进行演示，并通过数学模型外推至高刷新率配置下的指标，高刷新率硬件的真机验证属于未来工作。

\section{密码学安全伪随机数生成器}

掩模的随机性是本系统安全性的根基。若掩模由普通伪随机数生成器（如线性同余法）产生，攻击者有可能通过观测若干帧反推其内部状态，进而预测后续掩模。为此，本系统采用密码学安全伪随机数生成器（Cryptographically Secure Pseudo-Random Number Generator, CSPRNG）。

ChaCha20 是由 Bernstein 设计的流密码\textsuperscript{[10]}，具有高吞吐、抗时序攻击、无需查找表等优点，已被 TLS 1.3 等标准广泛采用。给定 256 位密钥和 96 位 nonce，ChaCha20 可生成统计上不可区分于真随机的密钥流，攻击者在不知密钥的情况下无法预测任意比特。本系统以 ChaCha20 密钥流驱动随机索引矩阵的生成。

为保证不同显示周期的掩模统计独立、防止固定图案被长曝光捕获，本系统采用基于 HMAC 的密钥派生函数（HMAC-based Key Derivation Function, HKDF）\textsuperscript{[11]}，以用户主密钥 $K$ 和周期序号 $t$ 派生周期子密钥 $K_t = \mathrm{HKDF}(K, \text{salt}=t)$。即使攻击者获取了某一周期的掩模，也无法借此推断其他周期的掩模。

\section{对抗样本技术}

对抗样本是指在原始输入上叠加精心构造的、人眼难以察觉的微小扰动，使机器学习模型产生错误输出的样本。本系统利用对抗样本技术，在子帧中注入对抗噪声，进一步降低单帧被识别的概率。

快速梯度符号法（FGSM）由 Goodfellow 等人提出\textsuperscript{[6]}，其核心思想是沿损失函数对输入梯度的符号方向施加扰动：

\begin{equation}
N = \epsilon \cdot \mathrm{sign}\big(\nabla_I\, \mathcal{L}(f_\theta(I), y_{\text{true}})\big)
\end{equation}

其中 $f_\theta$ 为目标识别模型，$\mathcal{L}$ 为损失函数，$\epsilon$ 为扰动预算（典型取 $8/255$），$\|N\|_\infty \le \epsilon$ 保证扰动有界、不可感知。投影梯度下降法（PGD）\textsuperscript{[7]}是 FGSM 的多步迭代版本，每步施加小幅扰动并投影回 $\epsilon$ 邻域，攻击能力更强，是当前公认的一阶对抗攻击基线。

本系统的关键创新在于：将基础对抗噪声 $N_{\text{base}}$ 分解为 $n$ 个满足 $\sum_{k=1}^{n} N_k = 0$ 的互补子噪声，使噪声在人眼积分后完全抵消（不影响观看），而在每个单帧内保持对抗强度（最大化机器识别错误率）。

\section{光学字符识别与目标检测技术}

光学字符识别（OCR）是将图像中的文字转换为可编辑文本的技术，是攻击者从屏幕截图中提取敏感信息的主要手段，因此也是本系统主要的防御对象与评测工具。本系统选用三种主流开源 OCR 引擎进行交叉验证：Tesseract 是历史悠久的传统 OCR 引擎，基于 LSTM 识别；EasyOCR 与 PaddleOCR 是基于深度学习的现代引擎，对自然场景文本鲁棒性更好。使用多引擎评测可避免对单一引擎过拟合，使防御效果结论更可信。

目标检测技术用于识别图像中物体的类别与位置。YOLO（You Only Look Once）系列\textsuperscript{[12]}是单阶段目标检测的代表，YOLOv8 在精度与速度上均有良好表现。本系统使用 YOLOv8 评估时间分片对"物体级场景信息"的防护效果，并与文本级 OCR 防护进行对比，以区分不同粒度信息的防护强度。

\section{本章小结}

本章介绍了本系统所依赖的关键技术：视觉暂留与人眼视觉积分特性提供了"人眼可见"的生理基础；时间分片与高刷新率显示给出了 $f_r \ge 60n$ 的工程约束；ChaCha20 CSPRNG 与 HKDF 保证了掩模的不可预测性与周期独立性；FGSM/PGD 对抗样本技术为单帧防护提供了算法支撑；OCR 与目标检测技术既是防御对象也是评测工具。这些技术的有机结合构成了下一章系统设计的基础。
